% =====================================================================
%  Chapter 22 — Selection: DSR, PBO, the Reality Check and the SPA
% =====================================================================

\chapter{Selection: DSR, PBO, the Reality Check and the SPA}\idx{Reality check}
\label{ch:selection}

This chapter builds the four tests that judge the book's own system. They
are not decorations and they are not applied to somebody else's strategy.
They are implemented here, from their definitions, so that the reader can
see exactly what is being claimed when the book says a result is
distinguishable from luck.

\section{Why one test is not enough}

Each of the four tests answers a different question, and none of them
answers all of it.

\begin{itemize}
\item The \textbf{Deflated Sharpe Ratio\idx{Deflated Sharpe ratio}} asks: given that I tried \(N\)
strategies, is my Sharpe ratio larger than the expected maximum Sharpe
ratio of \(N\) draws from the null?
\item The \textbf{Probability of backtest overfitting\idx{Probability of backtest overfitting}} asks: if I split my
record in half, how often does the strategy that wins on the first half
lose on the second?
\item The \textbf{Reality Check} asks: is the best member of my family
better than the benchmark by more than the family would manage by chance?
\item The \textbf{Superior Predictive Ability\idx{Superior predictive ability}} test asks the same
question, but weights the family members by their own performance rather
than assuming they are equally good.
\end{itemize}

A result that passes all four is stronger than a result that passes one.
A result that passes none is dead.

\section{The Deflated Sharpe Ratio}

Bailey and L\'opez de Prado's construction begins from the observation that
the expected maximum Sharpe ratio of \(N\) independent trials with
variance \(V\) across trials is
\begin{equation}
SR_0 = \sqrt{V}\left[(1-\gamma)\,Z^{-1}\!\left(1-\frac{1}{N}\right)
                     + \gamma\,Z^{-1}\!\left(1-\frac{1}{Ne}\right)\right],
\end{equation}
where \(Z^{-1}\) is the inverse standard normal cumulative distribution,
\(e\) is Euler's number, and \(\gamma\) is the Euler--Mascheroni constant,
\(0.5772156649\).

The Deflated Sharpe Ratio then asks whether the observed Sharpe ratio
exceeds \(SR_0\), correcting for the non-normality of the return
distribution:
\begin{equation}
DSR = Z\!\left[
  \frac{(\widehat{SR} - SR_0)\sqrt{T-1}}
       {\sqrt{1 - \gamma_3\widehat{SR}
              + \frac{\gamma_4 - 1}{4}\widehat{SR}^{\,2}}}
\right],
\end{equation}
where \(\gamma_3\) and \(\gamma_4\) are the skewness and kurtosis of the
returns, and \(Z\) is the standard normal cumulative distribution.

The correction term in the denominator is the skewness and kurtosis
adjustment: if the returns are negatively skewed, the observed Sharpe
ratio is flattered by the left tail, and the denominator grows. The book
computes \(\gamma_3\) and \(\gamma_4\) from the actual return series.

There is a second and subtler correction. The variance \(V\) in \(SR_0\)
is the variance of the Sharpe ratios \emph{across the trial family}, not
the variance of the returns. If the family members all produce nearly
identical portfolios --- which, as Chapter 33 will show, they nearly do ---
then \(V\) is tiny and \(SR_0\) is tiny, and the test is correspondingly
weak. That is correct behaviour, and it is why the book also computes the
implied number of independent trials and repeats the test with it.

\section{CSCV and the probability of backtest overfitting}

The probability of backtest overfitting is built on a combinatorial
construction that is worth understanding, because it is elegant and it is
the test that most directly answers the question a reader will ask: does
this strategy's edge survive being cut in half?

The construction is combinatorially symmetric cross-validation\idx{Combinatorially symmetric cross-validation} --- CSCV for
short --- and the symmetry in the name is the point of the construction:
each split is a cross-validation\idx{Cross-validation} in which the two halves exchange roles.

Take the \(T\) months of the strategy family's returns. Split them into
\(S\) contiguous sub-periods of equal length, in every possible way: the
number of splits is the binomial coefficient
\(\binom{T}{S/2}\), and the book uses eight sub-periods, so the splits
are the \(\binom{T}{4}\) ways of choosing four months out of eight.

For each split, the first half is the \emph{in-sample} half and the
second is the \emph{out-of-sample} half. For each split, compute the
Sharpe ratio of every family member on each half, rank the members by
Sharpe ratio on the first half, and record the rank of the member that
was best on the first half in the second half's ranking. The logit of the
normalised rank is
\begin{equation}
\lambda_c = \log\frac{\omega_c}{1-\omega_c},
\qquad
\omega_c = \frac{r^c_{n^*}}{N+1},
\end{equation}
where \(r^c_{n^*}\) is the out-of-sample rank of the in-sample winner and
\(N\) is the number of trials. The probability of backtest overfitting is
then the probability that \(\lambda < 0\), i.e. the probability that the
in-sample winner is out of sample in the bottom half:
\begin{equation}
PBO = \phi = \int_{-\infty}^{0} f(\lambda)\,d\lambda .
\end{equation}

The interpretation is devastating and is the reason this test is the one
a reader should insist on. A PBO of 0.5 means that a coin flip is as good
as the strategy for picking the configuration that will work. A PBO
above 0.5 means the selection procedure is actively harmful.

\section{The reality check}

White's reality check asks whether the best member of a family beats a
benchmark by more than the family would by chance. The test statistic is
the maximum, over the family, of the difference in mean performance
between the member and the benchmark, standardised:
\begin{equation}
V = \max_k \left(
  \frac{\bar{r}_{k} - \bar{r}_{b}}
       {\hat{\sigma}_{k}/\sqrt{T}}
  \right),
\qquad
d_k(t) = r_k(t) - r_b(t),
\end{equation}
with the null distribution obtained by bootstrapping the centred
differences \(d_k(t) - \bar{d}_k\).

The bootstrap\idx{Bootstrap} is the stationary bootstrap, in which each observation is
either drawn from the sample or from its immediate successor, with the
probability of drawing the successor fixed at a small value. The
stationary bootstrap is used rather than the simple bootstrap because the
returns are serially correlated, and the simple bootstrap would destroy
that dependence and produce a null distribution that is too narrow.

The reality check is the most conservative of the four tests, because it
does not assume the family members are equal and because it is testing the
\emph{maximum}, which is the statistic that a researcher would have
reported.

\section{Hansen's superior predictive ability test}

The SPA test corrects the reality check's most wasteful conservatism. The
reality check's null distribution is the distribution of the maximum over
the family, which is wide when the family contains both good and bad
members; a family with one good member among many mediocre ones will fail
the test not because the good member is not better, but because the
maximum of the null is wide.

Hansen reweights. The test statistic uses the studentised mean differences
with weights that are proportional to each member's own performance:
\begin{equation}
t_{k,t} = \frac{d_{k,t}}{\hat{\sigma}_k},
\qquad
\bar{t}_k = \frac{1}{T}\sum_{t=1}^{T} t_{k,t},
\qquad
w_k = \frac{\max(0, \bar{t}_k)}{\sum_{j}\max(0, \bar{t}_j)},
\end{equation}
and the test statistic is the studentised sum
\begin{equation}
t^* = \max_k \frac{\sqrt{T}\,\bar{t}_k}
                           {\sqrt{1 - kT/n}},
\end{equation}
with the same stationary bootstrap for the null distribution.

A member with a negative mean difference receives zero weight and does
not contribute to the width of the null. A family of one good member and
twenty bad ones therefore faces a much narrower null than under the
reality check, and the test has real power.

\section{How the book applies them}

All four tests are applied to the book's own family, against the
benchmark, with the following discipline.

The family is the one that was actually run: the model configurations,
the ablations, the sizing variants, the ensemble half-life variants and
the alternative strategic policies. Its size is published.

The tests are run against the benchmark on the same months as the system,
with the same alignment\idx{Alignment}, so that the two halves are directly comparable.

The implied number of independent trials is computed from the family's own
correlation matrix and substituted into the Deflated Sharpe Ratio, and
the result is reported both ways.

And the result is reported whatever it is. Chapter 33 is where the
numbers appear, and they are not flattering.

\begin{keytakeaways}
\item Four tests, four questions: is the Sharpe ratio above the expected
maximum for \(N\) trials; does the in-sample winner survive being cut in
half; does the best family member beat the benchmark by more than chance;
and the same question with the family weighted by its own performance.
\item The Deflated Sharpe Ratio deflates by the expected maximum Sharpe
ratio of the family and corrects for skewness and kurtosis.
\item The reality check and the SPA test differ in how they weight the
family, and the book's family deliberately contains members that the SPA
test's reweighting will discard.
\item All four are applied to the book's own family, with its size
published.
\end{keytakeaways}
